\documentclass[french]{article}

\usepackage[utf8]{inputenc}
\usepackage[T1]{fontenc}
\usepackage{lmodern}
\usepackage{geometry}
\usepackage{graphicx}
\usepackage{babel}
\usepackage{amsmath}
\usepackage{amsthm}
\usepackage{amssymb}
\usepackage{hyperref}
\usepackage{stmaryrd}
\usepackage{enumitem}
\usepackage{tcolorbox}
\usepackage{dsfont}
\usepackage{multirow}
\usepackage{etoc}
\usepackage{titlesec}
\usepackage{esint}
\usepackage{cancel}
\usepackage{mathdots}


\setlist[itemize]{label=$\bullet$}


\renewcommand{\P}{\mathbb{P}}
\newcommand{\N}{\mathbb{N}}
\newcommand{\Z}{\mathbb{Z}}
\newcommand{\Q}{\mathbb{Q}}
\newcommand{\R}{\mathbb{R}}
\newcommand{\C}{\mathbb{C}}
\newcommand{\K}{\mathbb{K}}
\renewcommand{\L}{\mathbb{L}}
\newcommand{\U}{\mathbb{U}}
\newcommand{\st}{^*}
\newcommand{\tq}{\middle|}
\newcommand{\inv}{^{-1}}
\newcommand{\E}{\mathbb{E}}
\newcommand{\F}{\mathbb{F}}
\newcommand{\V}{\mathbb{V}}
\renewcommand{\ker}{\text{Ker}}
\newcommand{\im}{\text{Im}}
\newcommand{\card}{\text{Card}}
\newcommand{\Tr}{\text{Tr}}
\newcommand{\id}{\text{id}}

\theoremstyle{definition}
\newtheorem{ex}{Exercice}
\newtheorem{dfn}{Définition}

\theoremstyle{theorem}
\newtheorem*{ind}{Indication}
\newtheorem{thm}{Théorème}
\newtheorem{cor}{Corollaire}

\theoremstyle{remark}
\newtheorem*{rmq}{Remarque}

\title{Entiers algébriques}
\begin{document}
\maketitle
\section{Cas des corps $\Q\subset\K\subset\C$}
\begin{dfn}[Entier algébrique]
	On dit qu'un nombre complexe est un entier algébrique lorsqu'il existe $P\in\Z[X]$ unitaire tel que $P(x)=0$. On notera $\mathcal{O}_\C$ l'ensemble des nombres complexes qui sont des entiers algébriques.
\end{dfn}
\begin{dfn}
	Si $\K$ est un corps tel que $\Q\subset\K\subset\C$, on définit $\mathcal{O}_\K=\K\cap\mathcal{O}_\C$.
\end{dfn}

\begin{thm}
	L'ensemble $\mathcal{O}_\C$ est un sous-anneau de $\C$.
\end{thm}
On en déduit immédiatement le corollaire suivant.
\begin{cor}
	Si $\K$ est un corps tel que $\Q\subset\K\subset\C$, alors $\mathcal{O}_\K$ est un sous-anneau de $\K$.
\end{cor}
Nous allons voir deux preuves différentes de ce théorème.
\subsection{Preuve par les modules de type fini}
\begin{proof}
	Soit $(x,y)\in\mathcal{O}_\C^2$. On écrit : $$\begin{cases}
		x^m=a_{m-1}x^{m-1}+\cdots+a_0 \\
		y^n=b_{n-1}x^{n-1}+\cdots+b_0 \\
	\end{cases}$$ avec les $a_i$ et les $b_j$ dans $\Z$. Par une récurrence immédiate, on en déduit que pour tout entier $k$ supérieur ou égal à $m$, $x^k$ est une combinaison linéaire à coefficients entiers de $x$. Il en va de même pour $y^l$ pour $l\geqslant n$. Alors, l'anneau $\Z[x,y]$ des $\Z$-combinaisons linéaires des $x^i y^j$ est engendré, en tant que groupe additif, par la famille des $(x^iy^j)$ pour $i\in\llbracket0,m-1\rrbracket$ et $j\in\llbracket0,n-1\rrbracket$. On pose $t=mn$ et on note $(u_k)_{1\leqslant k\leqslant t}$ une telle famille. On se donne $z\in\Z[x,y]$. Soit $k\in\llbracket1,t\rrbracket$. Comme $\Z[x,y]$ est un anneau, on sait que $zu_k\in\Z[x,y]$, et on écrit alors : $$zu_k=\sum_{l=1}^{t}\alpha_{l,k}u_l$$ avec les $\alpha_{l,k}$ dans $\Z$. On définit : $$M=\begin{pmatrix}
	z-\alpha_{1,1}&-\alpha_{1,2}&\cdots&-\alpha_{1,t}\\
	-\alpha_{2,1}&\ddots&&\vdots\\
	\vdots&&\ddots&\vdots\\
	-\alpha_{t,1}&\cdots&\cdots&z-\alpha_{t,t}
	\end{pmatrix}\ \ \text{et}\ \ U=\begin{pmatrix}
	u_1\\
	\vdots\\ 
	\vdots\\
	u_t
	\end{pmatrix}$$ On a alors, au vu des définitions précédentes, $MU=0$. En multipliant par la transposée de la comatrice de $M$ à gauche, on obtient : $$\det(M)U=0$$ Pour $k=1$, comme on a par exemple pris $u_k=1$, on obtient alors $\det(M)=0$. Or, $\det(M)$ est un polynôme en $z$, unitaire et à coefficients entiers. On en déduit que $z\in\mathcal{O}_\C$. En particulier, comme $x+y\in\Z[x,y]$ et $xy\in\Z[x,y]$, on en déduit que $\mathcal{O}_\C$ est un sous-anneau de $\C$.
\end{proof}

\subsection{Preuve par le résultant}
On se donne un anneau factoriel $A$. Pour $P$ et $Q$ deux polynômes de $A[X]$, on rappelle la notation $\text{Res}(P,Q)$ pour le résultant de $P$ et $Q$. On rappelle que $\text{Res}(P,Q)\in A$ et il vérifie : $$\text{Res}(P,Q)\ne 0\iff P\wedge Q=1$$
\begin{proof}
	On suppose que $(x,y)\in\mathcal{O}_\C^2$ avec $P$ et $Q$ unitaires à coefficients dans $\Z$ tels que $P(x)=Q(y)=0$. On pose $A=\Z[T]$ et $z=x+y$. Alors, $Q(T-X)\in A[X]$ donc $$R=\text{Res}(P(X),Q(T-X))\in A=\Z[T]$$ Comme $P$ et $Q$ sont unitaires, on peut obtenir avec la définition de $R$ que $R$ est unitaire. Alors : $$R(z)=\text{Res}(P(X),Q(z-X))$$ puisque l'évaluation est un morphisme. Or, $P(x)=0$ et $Q(z-x)=Q(y)=0$ donc $P$ et $Q(z-X)$ ne sont pas premiers entre eux. Donc par propriété du résultant, $R(z)=0$ si bien que $x+y\in\mathcal{O}_\C$. On peut faire de même pour $xy$ en considérant $$\text{Res}(P(X),X^{\deg(Q)}Q(TX\inv))$$ 
\end{proof}

\section{Cas des corps $\Q\subset\K$ de dimension finie sur $\Q$}
On se donne $\K$ un sur-corps de $\Q$ font la dimension en tant que $\Q$-espace vectoriel est finie. On note $\mathcal{O}_\K$ l'ensemble des éléments de $\K$ annulés par un polynôme unitaire à coefficients dans $\Z$.
\begin{dfn}
	Si $z\in\mathcal{O}_\K$, on appelle polynôme minimal de $z$ sur $\Q$ et on note $\mu_z$ le générateur unitaire du noyau de l'application linéaire : $$\begin{array}{lrcl}
		m_z:&\Q[X]&\longrightarrow&\K \\
		&Q&\longmapsto&Q(z)
	\end{array}$$
\end{dfn}
\begin{thm}
	Soit $z\in\mathcal{O}_\K$ et $\mu_z$ le polynôme minimal de $z$ sur $\Q$. Alors $\mu_z$ appartient à $\Z[X]$.
\end{thm}
\begin{proof}
	Fixons $P\in\Z[X]$ unitaire qui annule $z$. On peut alors écrire $P=\mu_z Q$ avec $Q\in\Q[X]$. On écrit $$\mu_z=\frac{a}{b}R\ \ \text{et}\ \ Q=\frac{c}{d}S$$ avec $(a,b,c,d)\in\Z^4$ et $c(R)=c(S)=1$ où $c(A)$ désigne le contenu du polynôme $A$. Ainsi : $$acRS=bdP$$ En passant aux contenus et par propriété du contenu, on a : $ac=bd$ puisque $c(P)=1$ car $P$ est unitaire. Donc $P=RS$. Quitte à multiplier par $-1$, $R$ et $S$ sont unitaires. Ainsi, $a=b$ en regardant le coefficient dominant de $\mu_z$, si bien que $\mu_z=R\in\Z[X]$.
\end{proof}
\begin{cor}
	On a $\mathcal{O}_\Q=\Z$.
\end{cor}
\begin{rmq}
	Lorsque l'anneau des entiers sur le corps des fractions d'une anneau de départ est égal à cet anneau de départ, on dit que l'anneau est intégralement clos. Ce théorème stipule donc que $\Z$ est intégralement clos. Dé façon plus générale, on a en fait montré que tout anneau factoriel est intégralement clos. La réciproque est fausse en général.
\end{rmq}
\begin{proof}
	L'inclusion réciproque est triviale car si $z\in\Z$, $X-z$ annule $z$ et est unitaire à coefficients entiers. Dans le sens direct, soit $z\in\mathcal{O}_\Q$. On a $X-z=\mu_z$. Mais $\mu_z\in\Z[X]$ d'après le théorème précédent.  Donc $z\in\Z$.
\end{proof}
\begin{rmq}[Un anneau ni factoriel, ni intégralement clos]
	Considérons $A=\C[X^2,X^3]$ le sous-anneau de $\C[X]$ engendré par $X^2$ et $X^3$. On a donc $A=\text{Vect}(1,X^2,X^3,\ldots)$. Alors $X=X^3/X^2\in\text{Frac}(A)$ mais $X\notin A$. Et $X$ est entier sur $A$ car $P=T^2-X^2\in A[T]$ annule $X$ et est unitaire.
\end{rmq}
\end{document}
